\uuid{sLAk}
\titre{Étude de fonction}
\auteur{bodin ; hanani ; mzari}
\organisation{exo7}
\licence{CC BY-NC-SA 4.0}
\datecreate{2020-01-20}
\niveau{L1}
\module{Analyse}
\chapitre{Fonctions circulaires et hyperboliques inverses}
\sousChapitre{Autre}
\difficulte{1}
\theme{Fonctions usuelles, étude de fonction}
\exo7qcmid{lille-1.fonctions-usuelles.010}

\contenu{
\begin{question}
\qclassification{126.00}
\qtags{motcle=étude de fonction}
Soit $f(x)=\frac{\ln x}{x}$.  Quelles sont les assertions vraies ?

\begin{answers}
    \good{$f$ est définie sur $]0,+\infty[$.}

    \bad{$f$ est croissante sur $]0,+\infty[$.}

    \good{$f$ est une bijection de $]0,e]$ dans $]-\infty, \frac{1}{e}]$.}

    \good{$f$ est une bijection de $[e,+\infty[$ dans $]0, \frac{1}{e}]$.}
\end{answers}
\begin{explanations}
$f$ est définie sur $]0,+\infty[$. En étudiant les variations de $f$, $f$ est strictement croissante sur $]0,e]$ et strictement décroissante sur $[e,+\infty[$. D'autre part, $f(]0,e])=]-\infty, \frac{1}{e}]$ et $f([e,+\infty[)=]0, \frac{1}{e}] $. On déduit que $f$ établit une bijection de $]0,e]$ dans $]-\infty, \frac{1}{e}]$ et de $[e,+\infty[$ dans $]0, \frac{1}{e}]$.
\end{explanations}
\end{question}
}
